\section{Related Work}
\label{sec:related_work}

\paragraph{Attacks on LLM agents and multi-agent systems.}
Attacks on LLM agents target the components that extend a base model into an
agent. At the single-agent level, direct and indirect prompt
injection~\citep{perez2022ignore,greshake2023indirect} hijack an agent through its
instructions or retrieved content, tool attacks abuse external
APIs~\citep{zhan2024injecagent}, and memory poisoning corrupts an agent's stored
knowledge~\citep{chen2024agentpoison}. In a multi-agent system these effects are
amplified rather than contained: a single compromised agent propagates toxic
content through communication and can paralyze the whole
network~\citep{gu2024agentsmith,yu2024netsafe,andriushchenko2024agentharm}.
Across these settings, however, the attacker is typically \emph{continuous}---it
acts maliciously from the outset and stays consistent across rounds, producing a
strong and persistent signal. The temporal profile of an attack, namely whether
its malicious signal stays constant or grows over time, has not been used as an
evasion lever. We fill this gap with an escalation attack whose signal is
deliberately weak in early turns and intensifies later, specifically to defeat
detectors that aggregate behavior over time.

\paragraph{Safeguarding agents and multi-agent systems.}
Defenses against agent attacks fall into single-agent guardrails and
topology-aware MAS defenses. Guardrail models such as
LlamaGuard~\citep{inan2023llamaguard} and WildGuard~\citep{han2024wildguard}
classify the inputs and outputs of an individual agent, but they are
topology-unaware and cannot reason about how malicious content propagates between
agents. Recent MAS defenses instead model the system as a graph and reason over
communication: NetSafe~\citep{yu2024netsafe} characterizes the topological safety
of multi-agent networks, while G-Safeguard~\citep{wang2025gsafeguard}, the
state-of-the-art and our primary baseline, detects attacked agents with a GNN over
a multi-agent utterance graph and remediates by pruning their edges. Their key
limitation, central to our work, is temporal: they compress each agent's
multi-turn dialogue into a single embedding---by mean-pooling or by keeping the
last turn---before message passing, discarding how an agent behaves over time. We
instead preserve cross-turn dynamics through a temporal edge encoder, which
exposes escalation while remaining effective against continuous attacks.

\paragraph{Multi-agent systems as graphs.}
A growing body of work represents MAS as graphs to organize or optimize
inter-agent communication. AutoGen~\citep{wu2023autogen} and
DyLAN~\citep{liu2023dylan} structure agent collaboration through predefined or
dynamic communication graphs, while GPTSwarm~\citep{zhuge2024gptswarm} and
AgentPrune~\citep{zhang2024agentprune} optimize graph topology---including by
pruning redundant edges---for performance and token economy. These methods use
the graph view for \emph{utility and efficiency} and operate on static or
aggregated message representations. We instead adopt the graph view for
\emph{security}, and, unlike prior graph-based detectors, equip the GNN with
temporal-aware edge features so that detection depends on how communication
evolves across rounds rather than on its average.
